\section{Reconstructing Human Control}
\label{sec:control}

Once AI is understood as an automatic instrumentality, the control question changes. The inquiry is not ``who selected this token?'' It is ``who constructed, supplied, authorized, observed, and could alter the deployment that produced the legally relevant effect?'' This shift matters because a model's internal computation is only partially observable, while the surrounding control structure is extensively documented in ordinary organizations.

This Part develops a narrow translation device: the \Cascade. It is adapted from systems-theoretic process analysis to represent an AI deployment as a set of controllers, control actions, process models, feedback channels, and stop capacities. Its function is evidentiary and organizational. It does not announce a new standard of care or calculate legal fault. It tells courts and parties where facts relevant to familiar doctrine are likely to be found.

\subsection{Partial Observability Is Not a Responsibility Vacuum}

Legal responsibility has never depended on complete access to an actor's internal state or to every physical step in a causal process. Intent and knowledge are commonly inferred from conduct and circumstances. Defect can be established circumstantially without identifying a specific internal failure. Organizational fault can lie in staffing, supervision, information systems, and resource allocation rather than a final employee's mental state. AI makes reconstruction difficult, but difficulty is not juridical absence.\lrfootnote{See, e.g., \casecite[444--46]{usGypsum1978} (explaining that intent ordinarily must be inferred from acts and surrounding circumstances); \casecite[41--43]{speller2003} (permitting circumstantial proof of product defect where the evidence excludes other causes not attributable to the defendant); \bluecite[\S\S~7.03(1), 7.05]{restatementAgency2006} (separating a principal's direct negligence from vicarious liability).}

Three forms of opacity should be separated. \term{Computational opacity} concerns the inability to give a compact human explanation of how internal model states produced a particular output. \term{system opacity} concerns hidden interactions among the model, prompts, retrieval, tools, filters, and software. \term{organizational opacity} concerns who approved the use, possessed risk information, controlled resources, or could stop operation. Only the first may be inherent to a model architecture. The second and third are often matters of design and recordkeeping.

Treating all three as ``the black box'' gives the least tractable fact too much power. A court may be unable to reconstruct why a model ranked one token above another yet able to determine that a provider selected the model, a deployer gave it patient records, an interface suppressed uncertainty, a safety team reported a known failure mode, management rejected a delay, and an operator lacked time to override. Those facts can establish or negate elements of an ordinary claim. Interpretability at the neuron level is not a condition precedent to responsibility at the system level.

The Hangzhou chatbot dispute makes the distinction visible. No employee contemporaneously decided to offer 100,000 yuan. But the provider controlled the service, issued warnings, selected technical precautions, and maintained the system through which the output appeared. The court could therefore reject contractual manifestation and still evaluate the provider's conduct under a fault-based duty.\lrfootnote{See \casecite{hangzhouChatbot2026}.} A theory that recognizes only immediate linguistic control would miss both halves of that analysis.

\subsection{The Cascade as a Deployment-Specific STPA Map}

Systems safety supplies a disciplined vocabulary for looking beyond a failed component. In STAMP and STPA, safety is treated as a control problem: losses can arise because constraints are absent, a controller supplies an unsafe action, feedback is inadequate, or a controller's process model does not match the system's actual state.\lrfootnote{See \bookcite{levesonEngineering2011}; \bluecite{levesonThomasSTPA2018}.} The method is particularly useful when software, people, management, and outside conditions interact and no single broken part explains the outcome.

Recent work has translated STPA to AI development and operation. Shalaleh Rismani, Roel Dobbe, and A. Jung Moon frame AI harm as a process-level hazard distributed across organizational roles rather than a defect discoverable only inside a model. A related seven-author study led by Rismani interviewed thirty industry practitioners and found both promise and nontrivial adoption costs in importing safety-engineering methods into responsible-ML practice.\lrfootnote{See \bluecite{rismaniSilos2026}; \artcite{rismaniPlaneCrashes2023}.} NIST's AI Risk Management Framework independently emphasizes lifecycle actors, documented responsibilities, executive accountability, testing, external feedback, third-party components, monitoring, and safe decommissioning.\lrfootnote{See \bluecite{nistAIRMF2023}.}

The \Cascade{} is an instantiation of those ideas for adjudication. It contains four functional levels:

\begin{enumerate}[leftmargin=2.2em]
\item \term{Authority and resource allocation}: the board, senior officials, budget holders, or public principals who approve objectives, risk tolerance, staffing, compute, data access, and release conditions.
\item \term{Provider and model team}: the actors who develop or select the model, document capability and limits, conduct evaluations, issue updates, and supply interfaces or weights.
\item \term{Deployment platform and model-harness operators}: the actors who choose system instructions, retrieval, memory, filters, routing, tools, credentials, user eligibility, monitoring, and escalation.
\item \term{User environment}: operators, users, affected persons, external systems, and physical or institutional conditions that supply inputs and experience effects.
\end{enumerate}

These are roles, not necessarily four firms. A vertically integrated provider can occupy the first three levels and operate the user-facing service, but affected persons and external conditions remain part of the environment rather than divisions of the provider. An API deployment can divide the controller roles among a model vendor, application developer, enterprise customer, and professional user. A public agency can be both resource allocator and deployer while a contractor supplies the model. The purpose of the map is to prevent a corporate name or contract label from substituting for actual function.

\begin{figure}[htbp]
\centering
\begin{tikzpicture}[
  node distance=0.82cm,
  box/.style={draw, rounded corners=2pt, align=center, text width=0.66\textwidth, minimum height=0.78cm, inner sep=5pt},
  down/.style={-{Latex[length=2mm]}, thick},
  up/.style={-{Latex[length=2mm]}, thick, dashed}
]
\node[box] (authority) {Authority and resources\\\footnotesize objectives, risk tolerance, budget, staffing, credentials, release approval};
\node[box, below=of authority] (provider) {Provider and model team\\\footnotesize model selection or development, evaluation, documentation, updates};
\node[box, below=of provider] (deployer) {Deployment platform and model harness\\\footnotesize instructions, retrieval, memory, tools, routing, safeguards, monitoring, override};
\node[box, below=of deployer] (environment) {User environment and affected domain\\\footnotesize operators, users, third-party systems, institutions, and experienced effects};

\draw[down] ([xshift=2.5cm]authority.south) -- node[right, font=\scriptsize, align=left] {constraints\\and resources} ([xshift=2.5cm]provider.north);
\draw[down] ([xshift=2.5cm]provider.south) -- node[right, font=\scriptsize, align=left] {control\\actions} ([xshift=2.5cm]deployer.north);
\draw[down] ([xshift=2.5cm]deployer.south) -- node[right, font=\scriptsize, align=left] {outputs and\\interventions} ([xshift=2.5cm]environment.north);

\draw[up] ([xshift=-2.5cm]environment.north) -- node[left, font=\scriptsize, align=right] {complaints, incidents,\\performance signals} ([xshift=-2.5cm]deployer.south);
\draw[up] ([xshift=-2.5cm]deployer.north) -- node[left, font=\scriptsize, align=right] {evaluations, exceptions,\\capacity limits} ([xshift=-2.5cm]provider.south);
\draw[up] ([xshift=-2.5cm]provider.north) -- node[left, font=\scriptsize, align=right] {risk reports, audit,\\resource requests} ([xshift=-2.5cm]authority.south);
\end{tikzpicture}
\caption{The deployment-control cascade. Solid arrows show authority, resources, and control actions moving toward operation. Dashed arrows show the feedback required to update each controller's process model. Override, rollback, credential revocation, and shutdown must be assigned at the level where they can be effective.}
\label{fig:cascade}
\end{figure}

For each level, the map asks six questions. What loss and hazard were reasonably within scope? Which control action could create or reduce the hazard? What did the controller believe about the system and environment? Which feedback could correct that belief? What authority and resources did the controller possess? Who could pause, roll back, constrain, or terminate operation? The answers are deployment specific. A generalized model card may inform them, but it cannot answer for a particular hospital, companion interface, or weapons-support workflow.

The stopping question deserves emphasis. ``Human control'' is often claimed whenever a person appears somewhere in a diagram. Effective intervention instead requires at least a detectable signal, sufficient information, time to act, authority over the relevant process, a usable control, and an institutional environment that permits its use. Automation can leave humans responsible for abnormal conditions while depriving them of the practice and information needed to respond---a problem described long before modern AI.\lrfootnote{See Lisanne Bainbridge, \artcite{bainbridgeIronies1983}.} A nominal reviewer who sees only the model's conclusion, has seconds to approve, and is penalized for deviation is a feedback endpoint, not a meaningful controller.

\subsection{Translating Engineering Variables into Legal Facts}

STPA does not determine duty, breach, causation, or remedy. Those are legal judgments supplied by the jurisdiction and cause of action. It does, however, identify observable variables that bear on them. The translation is direct enough to discipline discovery and expert proof:

\begin{longtable}{>{\raggedright\arraybackslash}p{0.23\textwidth} >{\raggedright\arraybackslash}p{0.31\textwidth} >{\raggedright\arraybackslash}p{0.36\textwidth}}
\caption{From Deployment Control to Doctrine}\label{tab:engineering-law}\\
\toprule
Control variable & Factual question & Potential doctrinal relevance \\
\midrule
\endfirsthead
\toprule
Control variable & Factual question & Potential doctrinal relevance \\
\midrule
\endhead
Controller & Who could set or enforce the relevant constraint? & Identity and scope of a possible duty holder; actual control rather than nominal title. \\
Control action & Who released, configured, routed, approved, updated, or withheld an intervention? & Conduct; authorization; attribution; breach; causal mechanism. \\
Process model & What did the actor believe about capability, users, environment, and safeguards? & Actual or constructive knowledge; foreseeability; reasonableness of a decision on the information available. \\
Feedback & Who received evaluations, incident reports, complaints, drift signals, overrides, and near misses? & Notice; continuing duty; opportunity to correct; state of mind where relevant. \\
Override, rollback, and stop authority & Who could reduce capability, revoke credentials, restore a prior version, or suspend operation, and on what signal? & Feasible precaution; ability to avoid or mitigate; causation; post-sale or post-deployment response. \\
Objectives, metrics, and resources & Who set engagement, speed, cost, accuracy, or safety targets and supplied the means to meet them? & Direct organizational negligence; incentive design; oversight; feasibility of compliance. \\
Constraint bypass & Who removed a safeguard, concealed a failure, exceeded a limit, or used unauthorized tools? & Individual conduct; comparative responsibility; superseding or intervening acts; internal recourse. \\
Logs and retention & Which versions, prompts, tool calls, approvals, alerts, and responses were preserved? & Proof; discovery proportionality; authentication; spoliation; rebuttal of causal inference. \\
\bottomrule
\end{longtable}

The middle column is deliberately factual. A person with shutdown credentials does not necessarily owe the plaintiff a duty; an actor who received an alert does not necessarily have legally sufficient notice; violation of an internal constraint is not automatically negligence per se. The mapping instead identifies evidence from which the operative doctrine can reason. It also exposes negative space. If no one can identify a risk owner, no channel carries complaints upward, or a deployed system cannot be rolled back, those are organizational design facts rather than mysteries attributable to the model.

The method improves causal analysis in two directions. A claimant can specify a safer control path: an age gate would have changed access; a confirmation step would have interrupted a purchase; a warning displayed at the crisis moment would have supplied different information; a stop rule triggered by repeated failures would have limited exposure. A defendant can identify a break in the path: a third party removed the safeguard, an independent deployer supplied undisclosed tools, a user deliberately defeated a reasonable constraint, or the proposed intervention could not have changed the outcome. ``AI unpredictability'' becomes a proposition to be tested against particular controls rather than an all-purpose conclusion.

\subsection{Stress Tests Across Mixed Deployments}

The cascade remains useful only if it survives fragmented architectures. Seven recurring deployments illustrate how responsibility can move without ever migrating to the machine.

\paragraph{Hosted, vertically integrated service.}
When one enterprise supplies the model, interface, memory, policies, monitoring, and updates, its practical control is concentrated. Vendors and contractors may contribute components, but the claimant-facing deployment is readily identifiable. Internal division among product, safety, legal, and operations belongs in the responsibility charter developed in Part V.

\paragraph{Model API and third-party application.}
The model provider can control base capability, access conditions, model updates, and some monitoring; the application deployer can control system prompts, retrieval, tools, user population, warnings, and business objectives. Contracts and technical documentation help identify those roles but do not conclusively allocate duties owed to strangers. Logs should show which layer produced, blocked, modified, or authorized the relevant effect.

\paragraph{Open weights and materially modified models.}
Publishing weights may end the original provider's practical ability to monitor or stop an independent downstream deployment. A modifier who removes safeguards, fine-tunes new behavior, adds credentials, and distributes an application occupies different control positions. The original release can still be examined under any duty applicable to that act; continuing responsibility should track continuing control and foreseeable use rather than attach indefinitely to provenance alone.

\paragraph{Anonymized intermediaries.}
A routing service, shell deployer, or undisclosed model source increases identification costs. It does not create an artificial responsible person. The service facing the user remains a logical starting point for process and discovery, while records can identify upstream components and allocation. A business choice not to retain essential provenance can itself become relevant under applicable recordkeeping and spoliation rules.

\paragraph{Government systems.}
Procurement officials may set objectives and contract conditions; a vendor may supply and update a model; a military, benefits, or law-enforcement component may control data and tools; an operator may implement recommendations. The public chain of command and the technical control path must be represented together. Classification as ``government AI'' does not collapse contractor and state functions, and a procurement clause cannot conclusively decide constitutional responsibility.

\paragraph{Human-in-the-loop systems.}
The location of a human checkpoint is the beginning of analysis. Meaningful review requires time, information independent of the recommendation, competence, authority to depart, a route to escalate, and protection against metrics that punish justified intervention. If those conditions are absent, the organization cannot point to the human icon as proof that control was transferred.

\paragraph{Tiered access and automated restriction.}
Public, professional, research, and government users may receive different models, tools, rate limits, or filters. Automated bans and escalation rules are themselves control actions. The cascade asks who selected the tiering criterion, what feedback detects error, and who can review a denial or expand access. Part VI considers when such differentiation requires transparency and review.

Across these cases, control is a vector rather than a switch. One actor may control training, another deployment context, a third credentials, and a fourth emergency shutdown. Where governing law supplies a duty, its scope and alleged breach should be tested against the risk and capability that actor actually governed. Part V turns this map into an external and internal responsibility structure.

The method has a final limit. STPA depends on the analyst's system boundary, selected losses, causal assumptions, and normative choice of constraints; its controller language can underdescribe cooperation, culture, and social power.\lrfootnote{See \bluecite{rismaniSilos2026}.} The cascade should therefore be contestable by parties and affected communities and revised as discovery changes the facts. It structures evidence. It does not decide duty, legal control, defect, fault, mens rea, proximate cause, or defense.
